% homework/hw04x/hw04x.tex
% Year 7 - Extra Homework: two-digit multiplication, and division.
%
% Goes out alongside Homework 4. As with hw03x, the framing is deliberate and
% nothing on the sheet says easy, basic, extra help, catch-up or revision. The
% cover states a fact instead, and it is a true one: the twelve squares stop at
% 12 x 12, and every square past that -- 13 x 13, 18 x 18, 24 x 24 -- IS a
% two-digit multiplication. Cubing is two of them. The check at the end of
% Homework 4 is 18 x 18 = 324 and then 324 x 18 = 5832. So this sheet is not
% beside Unit 1; it is the arithmetic Unit 1 finishes on.
%
% Two things are woven back into Homework 4 on purpose:
%   * E5 squares 13, 15, 16, 18, 24 and 25 -- 225 is the patio from the lesson,
%     576 is the worked example in Homework 4 B2, and 324 is the first half of
%     its Step 4 check.
%   * E10 is the four divisions Homework 4 actually needs (441 div 9, 1296 div
%     9, 1728 div 8, 5832 div 8), and the sheet says so. A student who does
%     them here has done them there.
%
% House style is shared: see homework/hw-style.tex. Method: the new-homework
% skill. Source is pure ASCII on purpose.
%
% Build:  pwsh .claude/skills/new-homework/scripts/build-hw.ps1 hw04x

\documentclass[12pt,a4paper]{article}

% -- Answer key switch -------------------------------------------------------
\newif\ifkey
\ifdefined\TEACHER\keytrue\else
  \keyfalse        % \keyfalse = student copy   |   \keytrue = teacher copy
\fi

% -- House style -------------------------------------------------------------
% homework/hw-style.tex
% Shared house style for every Year 7 homework packet. \input this from the
% packet file, right after \documentclass and the answer-key switch.
%
% Do not fork this file per packet. If a packet needs a different look, the
% question is whether every packet should have it.
%
% The choices here are deliberate and were arrived at by printing and looking:
%   * No solid dark banners. Every header is a light tint with a coloured spine
%     on the left, so colour reads clearly without flooding a school printer.
%   * No marks and no mark totals. These are practice, not exam papers.
%   * Lato at 12pt with open leading, plus mathastext so inline maths is Lato
%     too. Without mathastext every "-6 + 4" is Computer Modern serif and the
%     sheet looks half-typeset.
%   * Writing lines are spaced for Year 7 handwriting, not adult margin notes.
%
% Provides: \wlines \blank \tickbox \sep \qmk-free marks (none) \hwsection
% \qhead \expr, the workspace box, and the remember / wordhelp / activity /
% yourturn coloured boxes. Palette matches the lesson decks exactly.

% -- Packages ----------------------------------------------------------------
\usepackage[T1]{fontenc}
\usepackage[utf8]{inputenc}
\usepackage[default]{lato}
\usepackage[a4paper,top=1.7cm,bottom=1.7cm,left=2.0cm,right=2.0cm,headheight=15pt]{geometry}
\usepackage{amsmath,amssymb}
\usepackage{xcolor}
\usepackage[most]{tcolorbox}
\usepackage{enumitem}
\usepackage{tikz}
\usetikzlibrary{arrows.meta}
\usepackage{array}
\usepackage{tabularx}
\usepackage{fancyhdr}
\usepackage{multicol}
\usepackage{needspace}
\usepackage{graphicx}

% Make maths use Lato too. Without this, every "-6 + 4" on the page is set in
% Computer Modern serif and the sheet looks half-typeset. Must come last.
\usepackage[italic]{mathastext}

\linespread{1.06}\selectfont

% -- The Learner's Book palette (same hex values as the lesson decks) ---------
\definecolor{cbteal}{HTML}{0087A8}
\definecolor{cbpurple}{HTML}{5C2483}
\definecolor{cborange}{HTML}{C25E12}
\definecolor{cbgreen}{HTML}{4A8B23}
\definecolor{cbred}{HTML}{C8102E}
\definecolor{cbblue}{HTML}{1A5FA8}
\definecolor{rulegray}{HTML}{9AA5AD}

% -- Answer lines and blanks -------------------------------------------------
% \wlines{n} draws n ruled writing lines, spaced wide enough for Year 7
% handwriting rather than for an adult with a fine pen.
\makeatletter
\newcommand{\wlines}[1]{%
  \par\vspace{0.75em}%
  \@tempcnta=0\relax
  \@whilenum\@tempcnta<#1\do{%
    \hbox to \linewidth{\color{rulegray}\leaders\hrule height 0.5pt\hfill}%
    \vspace{1.5em}%
    \advance\@tempcnta by 1\relax}%
  \par\vspace{-0.8em}%
}
\makeatother

% \blank[width] is an inline gap to write a short answer in.
\newcommand{\blank}[1][2.6cm]{\underline{\hspace{#1}}}

% A boxed space for showing working.
\newtcolorbox{workspace}[1][3.4cm]{%
  colback=white, colframe=rulegray, boxrule=0.5pt, arc=5pt,
  height=#1, left=6pt, right=6pt, top=4pt, bottom=4pt,
  before skip=9pt, after skip=11pt}

% An empty tick box.
\newcommand{\tickbox}{\fbox{\rule{0pt}{1.15em}\rule{1.15em}{0pt}}}

% A middot separator.
\newcommand{\sep}{\,\textperiodcentered\,}

% -- Section and question headings -------------------------------------------
% Light tint plus a fat coloured spine on the left. No solid dark banner: the
% colour reads clearly but barely costs the printer anything.
% needspace stops a heading stranding alone at the foot of a page. Kept modest:
% too large and it pushes headings over, leaving a worse hole than it prevents.
\newcommand{\hwsection}[2][cbteal]{%
  \needspace{8\baselineskip}%
  \par\vspace{13pt}%
  \noindent
  \begin{tcolorbox}[enhanced, nobeforeafter, width=\textwidth,
      colback=#1!7, colframe=#1!7, arc=5pt,
      borderline west={5pt}{0pt}{#1},
      left=14pt, right=10pt, top=7pt, bottom=7pt]
    {\large\bfseries\color{#1}#2}
  \end{tcolorbox}%
  \par\vspace{9pt}}

% \qhead[colour]{A2}{Moving along the line} -- a question heading.
\newcommand{\qhead}[3][cbteal]{%
  \needspace{6\baselineskip}%
  \par\vspace{11pt}%
  \noindent{\large\bfseries\color{#1}#2}\quad{\large\bfseries #3}\par\vspace{4pt}}

% -- Coloured boxes ----------------------------------------------------------
% Shared look: rounded, light fill, coloured rule, coloured title text on a
% pale title strip rather than a dark one.
\tcbset{hwbox/.style={
  enhanced, breakable, arc=6pt, boxrule=1pt,
  left=11pt, right=11pt, top=8pt, bottom=8pt,
  fonttitle=\bfseries, attach boxed title to top left={xshift=12pt, yshift=-8pt},
  boxed title style={arc=4pt, boxrule=1pt, left=7pt, right=7pt, top=2pt, bottom=2pt},
  top=13pt}}

% Orange: things to remember. Matches the deck's "Write This Down".
\newtcolorbox{remember}[1][Remember from the lesson]{hwbox,
  colback=cborange!6, colframe=cborange, coltitle=cborange,
  colbacktitle=white, title={#1}}

% Blue: English help. The barrier in this class is the wording, not the maths.
\newtcolorbox{wordhelp}[1][Word help]{hwbox,
  colback=cbblue!6, colframe=cbblue, coltitle=cbblue,
  colbacktitle=white, title={#1}}

% Purple: an activity or instruction.
\newtcolorbox{activity}[1][How to do this homework]{hwbox,
  colback=cbpurple!6, colframe=cbpurple, coltitle=cbpurple,
  colbacktitle=white, title={#1}}

% Green: your turn / a friendly nudge.
\newtcolorbox{yourturn}[1][Your turn]{hwbox,
  colback=cbgreen!7, colframe=cbgreen, coltitle=cbgreen!75!black,
  colbacktitle=white, title={#1}}

% -- Shorthands --------------------------------------------------------------
\newcommand{\dC}{\,^{\circ}\mathrm{C}}          % use inside maths: $-8\dC$
\newcommand{\key}[1]{{\color{cborange}\textbf{#1}}}
\newcommand{\eng}[1]{{\color{cbblue}\textbf{#1}}}

\setlist[enumerate]{itemsep=8pt, topsep=5pt, leftmargin=*}
\setlist[itemize]{itemsep=4pt, topsep=4pt, leftmargin=*}
\renewcommand{\arraystretch}{1.45}

% -- An expression the student is asked to work from -------------------------
\newcommand{\expr}[1]{%
  \tcbox[on line, colback=cbgreen!10, colframe=cbgreen!55, boxrule=1pt, arc=4pt,
         left=9pt, right=9pt, top=4pt, bottom=4pt]{\large\bfseries $#1$}}

% -- Running head and foot ---------------------------------------------------
% The packet overrides \fancyhead[L] with its own title.
\pagestyle{fancy}
\fancyhf{}
\fancyhead[L]{\small\color{black!55}Year 7\sep Homework}
\fancyhead[R]{\small\color{black!55}Name: \blank[3.4cm]}
\fancyfoot[C]{\small\color{black!55}Page \thepage}
\renewcommand{\headrulewidth}{0pt}
\renewcommand{\headrule}{\vspace{2pt}\hbox to\headwidth{\color{cbteal!45}\leaders\hrule height 1.2pt\hfill}}



% -- Per-packet running head -------------------------------------------------
\fancyhead[L]{\small\color{black!55}Year 7\sep Extra Homework}

\begin{document}

% ===========================================================================
% COVER BLOCK
% ===========================================================================
\thispagestyle{fancy}

\begin{tcolorbox}[enhanced, colback=cbgreen!8, colframe=cbgreen!8, arc=7pt,
                  borderline west={6pt}{0pt}{cbgreen},
                  left=18pt, right=14pt, top=12pt, bottom=12pt]
  {\color{cbgreen!75!black}\bfseries YEAR 7\sep UNIT 1\sep EXTRA HOMEWORK}\\[6pt]
  {\color{cbgreen!70!black}\Huge\bfseries Two Rows, Then Add}\\[10pt]
  {\color{black!70}Two-digit multiplication\sep and the division that undoes it}
\end{tcolorbox}

\vspace{8pt}
\noindent
\begin{tabularx}{\textwidth}{@{}X X@{}}
  \textbf{Name:} \blank[5.6cm] & \textbf{Class:} \blank[5.6cm] \\
  \textbf{Date given:} \blank[4.6cm] & \textbf{Date due:} \blank[5.0cm] \\
\end{tabularx}

\vspace{6pt}
\begin{activity}[Why this sheet exists]
  The list of square numbers you copied stops at $12 \times 12$. Every square
  after that is a \key{two-digit multiplication}: $13 \times 13$, $18 \times 18$,
  $24 \times 24$. And a cube is \key{two} of them --- to check that
  $\sqrt[3]{5832} = 18$ you work out $18 \times 18 = 324$, and then
  $324 \times 18$.

  \medskip
  So this is not practice \textbf{beside} Unit 1. It is the arithmetic Unit 1
  \key{finishes on}, and the sheet ends on the exact divisions Homework 4 asks
  you for.

  \medskip
  \eng{No calculator}, and show your working --- working you can see is working
  you can check.
\end{activity}

% ===========================================================================
% PART 1 -- WARM UP
% ===========================================================================
\hwsection[cbgreen]{Part 1 \textperiodcentered\ The facts the method runs on}

\qhead[cbgreen]{E1}{Fill in the missing number}

\noindent Two-digit multiplication is these facts, done twice and added. Aim to
know them without counting.

\vspace{9pt}
\begin{multicols}{3}
\begin{enumerate}[label=(\alph*), itemsep=9pt]
  \item $6 \times 8 = \blank[1.3cm]$
  \item $7 \times 8 = \blank[1.3cm]$
  \item $9 \times 4 = \blank[1.3cm]$
  \item $7 \times 6 = \blank[1.3cm]$
  \item $8 \times 8 = \blank[1.3cm]$
  \item $9 \times 7 = \blank[1.3cm]$
  \item $6 \times 6 = \blank[1.3cm]$
  \item $9 \times 8 = \blank[1.3cm]$
  \item $7 \times 4 = \blank[1.3cm]$
  \item $9 \times 9 = \blank[1.3cm]$
  \item $8 \times 4 = \blank[1.3cm]$
  \item $7 \times 7 = \blank[1.3cm]$
  \item $12 \times 8 = \blank[1.3cm]$
  \item $11 \times 7 = \blank[1.3cm]$
  \item $12 \times 6 = \blank[1.3cm]$
\end{enumerate}
\end{multicols}

% -- E2 ---------------------------------------------------------------------
\qhead[cbgreen]{E2}{Multiply by a ten}

\begin{remember}[The second row is always a ten]
  To multiply by \key{20}, multiply by \key{2} and then \key{put a zero on the
  end}. $34 \times 2 = 68$, so $34 \times 20 = 680$. That is the whole trick, and
  it is what the second row of a two-digit multiplication is.
\end{remember}

\vspace{2pt}
\begin{multicols}{3}
\begin{enumerate}[label=(\alph*), itemsep=9pt]
  \item $23 \times 10 = \blank[1.7cm]$
  \item $23 \times 20 = \blank[1.7cm]$
  \item $41 \times 30 = \blank[1.7cm]$
  \item $17 \times 40 = \blank[1.7cm]$
  \item $56 \times 20 = \blank[1.7cm]$
  \item $38 \times 50 = \blank[1.7cm]$
\end{enumerate}
\end{multicols}

% ===========================================================================
% PART 2 -- TWO ROWS
% ===========================================================================
\hwsection[cbgreen]{Part 2 \textperiodcentered\ Two rows, then add}

\begin{remember}[The method, once, in full]
  \noindent
  \begin{minipage}[t]{0.30\linewidth}
    \vspace{2pt}
    {\renewcommand{\arraystretch}{1.05}%
    \begin{tabular}{r@{\ }r@{\ }r@{\ }r}
              &   & 4 & 7 \\
      $\times$&   & 2 & 6 \\[1pt]
      \hline
              & 2 & 8 & 2 \\
              & 9 & 4 & 0 \\[1pt]
      \hline
      1       & 2 & 2 & 2 \\[1pt]
      \hline
    \end{tabular}}
  \end{minipage}%
  \begin{minipage}[t]{0.70\linewidth}
    \vspace{2pt}
    \key{Row 1:} $47 \times 6 = 282$. \\
    \key{Row 2:} $47 \times 20 = 940$ --- multiply by 2, then put the zero on. \\
    \key{Add them:} $282 + 940 = 1222$. \\[4pt]
    \eng{If you only write one row, you have multiplied by 6 and forgotten the
    20.} Your answer will come out about ten times too small. \textbf{Two rows,
    every time.}
  \end{minipage}
\end{remember}

% -- E3 ---------------------------------------------------------------------
\qhead[cbgreen]{E3}{Warm up: two digits by one digit}

\noindent Set each one out properly and show your working in the box.

\vspace{4pt}
\noindent
\begin{tabularx}{\textwidth}{|X|X|X|X|}
  \hline
  $46 \times 7$ & $58 \times 6$ & $73 \times 8$ & $89 \times 4$ \\
  \rule[-2.6em]{0pt}{3.2em} & & & \\
  \hline
\end{tabularx}

% -- E4 ---------------------------------------------------------------------
\qhead[cbgreen]{E4}{Two digits by two digits}

\noindent \key{Two rows, then add.} Show both rows every time.

\vspace{4pt}
\noindent
\begin{tabularx}{\textwidth}{|X|X|X|}
  \hline
  $23 \times 14$ & $34 \times 16$ & $27 \times 25$ \\
  \rule[-3.5em]{0pt}{4.3em} & & \\
  \hline
\end{tabularx}

\vspace{9pt}
\noindent
\begin{tabularx}{\textwidth}{|X|X|X|}
  \hline
  $58 \times 17$ & $39 \times 41$ & $75 \times 36$ \\
  \rule[-3.5em]{0pt}{4.3em} & & \\
  \hline
\end{tabularx}

\vspace{9pt}
\noindent
\begin{tabularx}{\textwidth}{|X|X|X|}
  \hline
  $46 \times 32$ & $64 \times 23$ & $52 \times 48$ \\
  \rule[-3.5em]{0pt}{4.3em} & & \\
  \hline
\end{tabularx}

\vspace{9pt}
\noindent
\begin{tabularx}{\textwidth}{|X|X|X|}
  \hline
  $84 \times 26$ & $47 \times 53$ & $96 \times 15$ \\
  \rule[-3.5em]{0pt}{4.3em} & & \\
  \hline
\end{tabularx}

\vspace{10pt}
% The blank has to be short enough to sit on the same line as the question. At
% 6.0cm and again at 5.0cm this ran past the right margin.
\noindent\textbf{Look back.} Two of those twelve have the \textbf{same answer}.
Which two? \blank[4.4cm]

% -- E5 ---------------------------------------------------------------------
\qhead[cbgreen]{E5}{Squaring: the same number, twice}

\begin{remember}[Where your list of squares runs out]
  Your notebook has the squares up to $12 \times 12 = 144$. After that there is no
  list --- you \key{multiply}. $13^2$ means $13 \times 13$, and that is a
  two-digit multiplication like any other.
\end{remember}

\noindent Work each one out. \key{You will meet three of these answers again in
Homework 4.}

\vspace{4pt}
\noindent
\begin{tabularx}{\textwidth}{|X|X|X|}
  \hline
  $13 \times 13$ & $15 \times 15$ & $16 \times 16$ \\
  \rule[-3.5em]{0pt}{4.3em} & & \\
  \hline
\end{tabularx}

\vspace{9pt}
\noindent
\begin{tabularx}{\textwidth}{|X|X|X|}
  \hline
  $18 \times 18$ & $24 \times 24$ & $25 \times 25$ \\
  \rule[-3.5em]{0pt}{4.3em} & & \\
  \hline
\end{tabularx}

\vspace{10pt}
\noindent\textbf{Now use one.} Your answer to $18 \times 18$ is the first half of a
check in Homework 4. Finish it here --- this one is \textbf{three digits by two},
so it takes two rows like all the rest.

\vspace{4pt}
\noindent\hspace*{1em}$18 \times 18 = \blank[2.4cm]$, and then
$\blank[2.4cm] \times 18 = \blank[3.0cm]$

\begin{workspace}[3.2cm]\end{workspace}

% ===========================================================================
% PART 3 -- DIVISION
% ===========================================================================
\hwsection[cbgreen]{Part 3 \textperiodcentered\ Division, which undoes it}

\qhead[cbgreen]{E6}{Divide --- these all come out exactly}

\begin{remember}[A check that costs five seconds]
  When you finish a division, \key{multiply your answer back}. If
  $84 \div 6 = 14$, then $14 \times 6$ must come to 84. If it does not, you have
  just caught your own mistake.
\end{remember}

\vspace{2pt}
\begin{multicols}{2}
\begin{enumerate}[label=(\alph*), itemsep=12pt]
  \item $84 \div 6 = \blank[2.2cm]$
  \item $92 \div 4 = \blank[2.2cm]$
  \item $135 \div 5 = \blank[2.2cm]$
  \item $168 \div 7 = \blank[2.2cm]$
  \item $216 \div 8 = \blank[2.2cm]$
  \item $288 \div 12 = \blank[2.2cm]$
  \item $405 \div 9 = \blank[2.2cm]$
  \item $728 \div 8 = \blank[2.2cm]$
  \item $924 \div 7 = \blank[2.2cm]$
  \item $576 \div 4 = \blank[2.2cm]$
\end{enumerate}
\end{multicols}

% -- E7 ---------------------------------------------------------------------
\qhead[cbgreen]{E7}{Divide --- these do not}

\begin{wordhelp}[Word help --- the words that describe what is left]
  {\renewcommand{\arraystretch}{1.15}%
  \begin{tabularx}{\linewidth}{@{}l@{\hspace{1.4em}}X@{}}
    \eng{exactly}    & nothing is left over. $84 \div 6 = 14$ exactly. \\
    \eng{remainder}  & what is left over when it does not come out exactly. We
                       write $75 \div 8 = 9$ r 3. \\
    \eng{goes into}  & ``8 goes into 24 three times'' means $24 \div 8 = 3$. \\
    \eng{left over}  & the everyday English for a remainder. \\
  \end{tabularx}}
\end{wordhelp}

\noindent Write each answer as a whole number and a remainder, like this:
\quad $75 \div 8 = 9$ r 3.

\vspace{6pt}
\begin{multicols}{2}
\begin{enumerate}[label=(\alph*), itemsep=12pt]
  \item $100 \div 6 = \blank[2.8cm]$
  \item $131 \div 9 = \blank[2.8cm]$
  \item $250 \div 7 = \blank[2.8cm]$
  \item $194 \div 12 = \blank[2.8cm]$
  \item $365 \div 4 = \blank[2.8cm]$
  \item $500 \div 8 = \blank[2.8cm]$
  \item $425 \div 6 = \blank[2.8cm]$
  \item $800 \div 9 = \blank[2.8cm]$
\end{enumerate}
\end{multicols}

% -- E8 ---------------------------------------------------------------------
\qhead[cbgreen]{E8}{Find the missing number}

\begin{remember}[Multiplying and dividing undo each other]
  If $14 \times 13 = 182$, then $182 \div 14 = 13$ \key{and} $182 \div 13 = 14$.
  Every multiplication fact is also \key{two} division facts.
\end{remember}

\vspace{2pt}
\begin{multicols}{2}
\begin{enumerate}[label=(\alph*), itemsep=12pt]
  \item $14 \times \blank[1.6cm] = 182$
  \item $\blank[1.6cm] \times 16 = 384$
  \item $576 \div \blank[1.6cm] = 24$
  \item $\blank[1.6cm] \div 9 = 49$
  \item $25 \times \blank[1.6cm] = 625$
  \item $\blank[1.6cm] \times 18 = 324$
  \item $1296 \div \blank[1.6cm] = 144$
  \item $\blank[1.6cm] \div 8 = 216$
\end{enumerate}
\end{multicols}

% -- E9 ---------------------------------------------------------------------
% This one is here because the word problems need it: two of the three divide
% by 24, and a student who has never divided by a two-digit number will stall
% there and look like they cannot read the question.
\qhead[cbgreen]{E9}{Divide by a two-digit number}

\begin{remember}[Same method, bigger question]
  Dividing by 24 is the same as dividing by 6. You still ask \key{how many fit},
  you are just asking it about a bigger number. It helps to write the first few
  out first: $24, 48, 72, 96, 120, \ldots$

  \smallskip
  And the check is the same: \key{multiply your answer back}.
\end{remember}

\noindent These all come out exactly. Show your working in the box.

\vspace{4pt}
\noindent
\begin{tabularx}{\textwidth}{|X|X|}
  \hline
  $384 \div 16$ & $950 \div 25$ \\
  \rule[-3.2em]{0pt}{4.0em} & \\
  \hline
\end{tabularx}

\vspace{9pt}
\noindent
\begin{tabularx}{\textwidth}{|X|X|}
  \hline
  $672 \div 24$ & $585 \div 13$ \\
  \rule[-3.2em]{0pt}{4.0em} & \\
  \hline
\end{tabularx}

% -- E10 --------------------------------------------------------------------
\qhead[cbgreen]{E10}{Use it: the four divisions Homework 4 asks for}

\begin{remember}[Straight out of Homework 4]
  In Homework 4 you split a big number into two pieces you know the root of. The
  \key{splitting is a division}, and these are the four it needs. Do them here and
  you have done them there.
\end{remember}

\vspace{4pt}
\noindent
\begin{tabularx}{\textwidth}{|p{3.0cm}|X|p{5.2cm}|}
  \hline
  \textbf{The division} & \textbf{Answer} & \textbf{So the split is} \\
  \hline
  $441 \div 9$  \rule[-0.9em]{0pt}{2.45em} & & $441 = 9 \times \blank[1.6cm]$ \\ \hline
  $1296 \div 9$ \rule[-0.9em]{0pt}{2.45em} & & $1296 = 9 \times \blank[1.6cm]$ \\ \hline
  $1728 \div 8$ \rule[-0.9em]{0pt}{2.45em} & & $1728 = 8 \times \blank[1.6cm]$ \\ \hline
  $5832 \div 8$ \rule[-0.9em]{0pt}{2.45em} & & $5832 = 8 \times \blank[1.6cm]$ \\ \hline
\end{tabularx}

% -- E11 --------------------------------------------------------------------
\qhead[cbgreen]{E11}{Word problems}

\begin{wordhelp}[Word help --- which operation does the sentence want?]
  {\renewcommand{\arraystretch}{1.15}%
  \begin{tabularx}{\linewidth}{@{}l@{\hspace{1.4em}}X@{}}
    \eng{each} / \eng{every}          & usually multiply. 24 boxes with 36 in
                                        \textbf{each} is $24 \times 36$. \\
    \eng{share equally between}       & divide. \\
    \eng{how many are left over}      & the remainder. \\
    \eng{altogether} / \eng{in total} & the whole amount, once you have finished. \\
  \end{tabularx}}
\end{wordhelp}

\noindent Write the calculation first, then the answer.

\begin{enumerate}[label=\textbf{\arabic*.}, itemsep=11pt]

  \item \textbf{The exercise books.} Mr Bowen orders \textbf{24 packs} of exercise
        books. There are \textbf{36 books} in each pack. How many books does he
        have altogether?
        \begin{workspace}[3.0cm]\end{workspace}

  \item \textbf{The chairs.} There are \textbf{950 chairs} in the hall, set out in
        rows of \textbf{24}. How many full rows are there, and how many chairs are
        left over?
        \begin{workspace}[3.0cm]\end{workspace}

  \item \textbf{The wall.} Mr Bowen tiles a wall with \textbf{26 rows} of
        \textbf{34 tiles}. How many tiles does the wall take altogether?
        \begin{workspace}[3.0cm]\end{workspace}

  \item \textbf{The pencils.} Mr Bowen buys \textbf{27 boxes} of pencils with
        \textbf{36 pencils} in each box. He shares them equally between his
        \textbf{24 students}, and keeps whatever is left over for himself. How
        many pencils does each student get, and how many does Mr Bowen keep?
        % The two-step one, and the last thing on the sheet: the leftover space at
        % the foot of the last page belongs to the student, not to the margin.
        \begin{workspace}[5.2cm]\end{workspace}

\end{enumerate}

% ===========================================================================
% ANSWER KEY (teacher copy only)
% ===========================================================================
\ifkey
\newpage
\hwsection[cbred]{Answer Key --- Teacher Copy}

\linespread{1.0}\selectfont
\setlist[enumerate]{itemsep=2pt, topsep=3pt, leftmargin=*}
\setlist[itemize]{itemsep=2pt, topsep=3pt, leftmargin=*}

\qhead[cbred]{E1}{Fill in the missing number}
\noindent (a) 48 \quad (b) 56 \quad (c) 36 \quad (d) 42 \quad (e) 64 \quad (f) 63
\quad (g) 36 \quad (h) 72 \quad (i) 28 \quad (j) 81 \quad (k) 32 \quad (l) 49
\quad (m) 96 \quad (n) 77 \quad (o) 72

\smallskip
\noindent{\small Time this one rather than marking it. If a student can do all
fifteen in under two minutes, the tables are not what is slowing them down and
Part 2 is where to look instead. The three worth re-drilling if they go wrong are
(b) $7 \times 8$, (f) $9 \times 7$ and (h) $9 \times 8$ --- those three account for
most of the errors in a two-row multiplication, because they turn up in both rows.}

\qhead[cbred]{E2}{Multiply by a ten}
\noindent (a) 230 \quad (b) 460 \quad (c) 1230 \quad (d) 680 \quad (e) 1120
\quad (f) 1900

\smallskip
\noindent{\small This question is the second row of E4 taken out and practised on
its own, which is why it comes before the method rather than after it. The error to
look for is a missing zero --- 46 instead of 460 in (b). A student who makes it here
will make it in every question in E4, and it is much faster to catch on this line
than inside a four-line sum.}

\qhead[cbred]{E3}{Warm up: two digits by one digit}
\noindent $46 \times 7 = 322$ \quad $58 \times 6 = 348$ \quad $73 \times 8 = 584$
\quad $89 \times 4 = 356$

\smallskip
\noindent{\small The error here is carrying, not the tables. $89 \times 4$ carries
at both columns and is the one to check.}

\qhead[cbred]{E4}{Two digits by two digits}
\noindent
\begin{tabularx}{\linewidth}{@{}p{4.6cm}p{4.6cm}X@{}}
  $23 \times 14 = 322$  & $34 \times 16 = 544$  & $27 \times 25 = 675$ \\
  $58 \times 17 = 986$  & $39 \times 41 = 1599$ & $75 \times 36 = 2700$ \\
  $46 \times 32 = 1472$ & $64 \times 23 = 1472$ & $52 \times 48 = 2496$ \\
  $84 \times 26 = 2184$ & $47 \times 53 = 2491$ & $96 \times 15 = 1440$ \\
\end{tabularx}

\smallskip
\noindent \textbf{Look back:} $46 \times 32$ and $64 \times 23$ --- both come to
\textbf{1472}.

\smallskip
\noindent{\small Mark the \emph{working}, not just the answer. \textbf{One row of
working is the whole diagnosis}: it means the student multiplied by the units digit
and stopped, and the answer will be roughly a tenth of the right size --- 322
becomes 92, 1472 becomes 92. That is a method error and needs re-teaching. An answer
out by a small amount with two rows present is arithmetic, and needs a re-check, not
a re-teach. The \emph{Look back} pair is not decoration: 46, 32 and 64, 23 are the
same digits swapped, and it makes them read back their own twelve answers, which is
the check they never do voluntarily.}

\qhead[cbred]{E5}{Squaring: the same number, twice}
\noindent $13^2 = 169$ \quad $15^2 = 225$ \quad $16^2 = 256$ \quad
$18^2 = 324$ \quad $24^2 = 576$ \quad $25^2 = 625$

\smallskip
\noindent \textbf{Now use one:} $18 \times 18 = 324$, and $324 \times 18 = 5832$.

\smallskip
\noindent{\small Three of these are already in the class's week and it is worth
saying so out loud. \textbf{225} is the patio from the lesson --- 15 stones along
each side. \textbf{576} is the worked example in Homework 4 B2. \textbf{324} is the
first half of the check in Homework 4 B3, and the second half, $324 \times 18$, is
the only three-by-two multiplication on this sheet: it is meant to be hard, and a
student who reaches 5832 has proved a cube root by hand.}

\qhead[cbred]{E6}{Divide --- these all come out exactly}
\noindent (a) 14 \quad (b) 23 \quad (c) 27 \quad (d) 24 \quad (e) 27 \quad (f) 24
\quad (g) 45 \quad (h) 91 \quad (i) 132 \quad (j) 144

\smallskip
\noindent{\small (h) $728 \div 8 = 91$ is the one to point at afterwards: it is
exactly the test-for-8 working from Homework 4, where 1728 has to be shown divisible
by 8 before it can be split. Note that (c) and (e) both come to 27 and (d) and (f)
both come to 24 --- if a student notices, that is a good sign.}

\qhead[cbred]{E7}{Divide --- these do not}
\noindent (a) 16 r 4 \quad (b) 14 r 5 \quad (c) 35 r 5 \quad (d) 16 r 2 \quad
(e) 91 r 1 \quad (f) 62 r 4 \quad (g) 70 r 5 \quad (h) 88 r 8

\smallskip
\noindent{\small The diagnosis is in the remainder. A remainder \textbf{bigger than
the divisor} is a method error --- the student stopped dividing too early --- and
needs re-teaching. A remainder wrong by one or two is arithmetic and needs a
re-check. Only the first is worth class time.}

\qhead[cbred]{E8}{Find the missing number}
\noindent (a) 13 \quad (b) 24 \quad (c) 24 \quad (d) 441 \quad (e) 25 \quad
(f) 18 \quad (g) 9 \quad (h) 1728

\smallskip
\noindent{\small (d) and (h) are the hard ones, because the missing number is the
\emph{big} one and the student has to multiply rather than divide: $49 \times 9$ and
$216 \times 8$. A student who tries to divide there has not seen which number is
missing. Every answer in the second half is a number from Homework 4 --- 441, 324,
1296, 1728 --- so this question is worth going over before the packet is due rather
than after.}

\qhead[cbred]{E9}{Divide by a two-digit number}
\noindent $384 \div 16 = \mathbf{24}$ \quad $950 \div 25 = \mathbf{38}$ \quad
$672 \div 24 = \mathbf{28}$ \quad $585 \div 13 = \mathbf{45}$

\smallskip
\noindent{\small This question exists because of E11, where two of the four word
problems divide by 24. A student who has never divided by a two-digit number stalls
there and looks as though they could not read the question, so check this before
concluding it was the English. $585 \div 13$ is the hardest; a student who wrote out
13, 26, 39, 52 first has the right instinct.}

\qhead[cbred]{E10}{Use it: the four divisions Homework 4 asks for}
\noindent $441 \div 9 = \mathbf{49}$, so $441 = 9 \times 49$. \\
$1296 \div 9 = \mathbf{144}$, so $1296 = 9 \times 144$. \\
$1728 \div 8 = \mathbf{216}$, so $1728 = 8 \times 216$. \\
$5832 \div 8 = \mathbf{729}$, so $5832 = 8 \times 729$.

\smallskip
\noindent{\small This is the join between the two sheets and it is worth saying out
loud when you hand them out: the divisibility test picks the split, but it is a
division that carries it out, and that division is this. $5832 \div 8$ is the hardest
on the sheet and the one Homework 4 needs most. A student who arrives at Homework 4
with these four already done can spend the time on the roots instead of on the long
division, which is the point of handing the two out together.}

\qhead[cbred]{E11}{Word problems}
\begin{enumerate}[label=\arabic*., itemsep=2pt]
  \item $24 \times 36 = \mathbf{864}$ books.
  \item $950 \div 24 = 39$ r 14: \textbf{39 full rows}, with \textbf{14 chairs}
        left over.
  \item $26 \times 34 = \mathbf{884}$ tiles.
  \item $27 \times 36 = 972$ pencils. $972 \div 24 = 40$ r 12, so each student gets
        \textbf{40} and Mr Bowen keeps \textbf{12}.
\end{enumerate}
\noindent{\small Problem 4 is the only two-step question: multiply, then divide with
a remainder. The common failure is not arithmetic but stopping after step one and
answering 972 --- read the question back to them rather than correcting it. Problem 2
is the sentence test: ``39.58'' is a correct division answering a question nobody
asked, because you cannot have most of a row of chairs. Problems 1 and 3 are E4
wearing a story, so a student who can do E4 and not these has a reading gap, not an
arithmetic one.}

\fi
\end{document}
